\documentclass[11pt]{article}
\usepackage[margin=0.50in]{geometry}
\usepackage[T1]{fontenc}
\usepackage{array}
\usepackage{booktabs}
\usepackage{hyperref}
\usepackage{xcolor}

\definecolor{quietgreen}{HTML}{355E3B}
\definecolor{claimink}{HTML}{1F2421}

\hypersetup{
  colorlinks=true,
  urlcolor=quietgreen,
  linkcolor=quietgreen
}

\setlength{\parindent}{0pt}
\setlength{\parskip}{2.6pt}
\renewcommand{\arraystretch}{1.08}
\pagestyle{empty}

\begin{document}
\pagecolor{white}
\color{claimink}

{\Large\bfseries\color{claimink} ClaimReady: Small Claims, Filed Right}\\[-1pt]
{\color{quietgreen}\rule{\linewidth}{0.8pt}}\\[-3pt]
{\large\bfseries\color{quietgreen} For the invoice too small for a lawyer and too large to write off.}\\[-1pt]
ClaimReady helps NYC freelancers turn messy contracts, invoices, emails, and screenshots into a court-ready small-claims packet for unpaid service work.

{\color{quietgreen}\rule{\linewidth}{0.35pt}}\vspace{-3pt}

{\bfseries The user.} The concrete user is a New York City freelancer or solo service provider owed \(\$1{,}000\)--\(\$10{,}000\) by a NY-registered LLC or corporation: a designer, developer, photographer, writer, tutor, consultant, or marketer with a contract, invoice, emails, screenshots, or notes showing completed work and nonpayment. This is not a niche edge case. Upwork estimates that 64 million Americans freelanced in 2023, or 38\% of the U.S. workforce. A New York freelancer survey found that 62\% experienced nonpayment; among those who lost income, 51\% lost more than \(\$1{,}000\) and 22\% lost more than \(\$5{,}000\).

{\bfseries The problem.} Today, the user sends awkward follow-up emails, copies a generic demand-letter template, gives up, or tries to file in small claims alone. Hiring counsel rarely makes sense: Clio reports the average New York lawyer hourly rate at \(\$426\), which can consume the economics of a \(\$2{,}000\)--\(\$5{,}000\) claim. Small claims is designed for this gap, but filing still requires the user to identify the correct legal defendant, find a service address, choose venue, organize exhibits, calculate damages, and create forms the clerk can process. NYC Small Claims allows claims up to \(\$10{,}000\) with only a \(\$15\)--\(\$20\) filing fee, but the paperwork and confidence gap stop many people before they start.

{\bfseries The economics.} For one user-month, assume one case packet: about 120k input tokens for intake, evidence, retrieval, and tool context, plus 15k output tokens for extracted facts, agent outputs, and final drafting. With Gemini 2.5 Flash pricing around \(\$0.30\) per 1M input tokens and \(\$2.50\) per 1M output tokens, model cost is about \(\$0.07\) per case. Cloud Run, storage, PDF generation, logs, failed runs, payment fees, and light support add an estimated \(\$3\)--\(\$5\) of variable cost; the model is not the cost driver---infrastructure and Stripe fees are.

{\itshape Pricing ladder.} ClaimReady starts with simple flat pricing, then adds value capture once usage proves demand. \textbf{Phase 1:} \(\$29\) per packet to validate willingness-to-pay and gather the first thousand testimonials and defendant patterns. \textbf{Phase 2:} \(\$29\) plus 1--2\% of claim value at filing, charged upfront rather than as a contingency on recovery; trigger this after roughly 500 sustained cases/month. \textbf{Phase 3:} B2B distribution inside freelancer-facing platforms---Upwork, Fiverr, FreshBooks, QuickBooks Self-Employed---at a \(\$10\)--\(\$15\) wholesale price, with the partner paying acquisition.

\vspace{-1pt}
\begin{center}
\begin{tabular}{>{\raggedright\arraybackslash\bfseries}p{1.42in}>{\raggedright\arraybackslash}p{1.55in}>{\raggedright\arraybackslash}p{1.85in}>{\raggedright\arraybackslash}p{1.85in}}
\toprule
Item & Phase 1 (Launch) & Phase 2 (Scale) & Phase 3 (Distribution) \\
\midrule
Price & \(\$29\) / packet & \(\$29\) + 1.5\% of claim value & \(\$10\)--\(\$15\) wholesale \\
Avg revenue / case & \(\$29\) & \(\sim\$81\) (avg claim \(\$3.5\)k) & \(\sim\$12\) \\
Variable cost & \(\$3\)--\(\$5\) & \(\$3\)--\(\$5\) & \(\$3\)--\(\$5\) \\
Contribution / case & \(\sim\$25\) (\(\sim\)85\%) & \(\sim\$76\) (\(\sim\)94\%) & \(\sim\$8\) (\(\sim\)67\%) \\
CAC payback & \(<1.5\times\) / case & \(\sim\)3\(\times\) / case & Partner-paid CAC \\
\bottomrule
\end{tabular}
\end{center}
\vspace{-1pt}

{\itshape CAC and payback.} Blended CAC at launch is estimated at \(\$15\)--\(\$30\): SEO on the ``how do I sue an LLC in NY'' long-tail, Reddit communities for freelancers and personal finance, freelancer Twitter, and partnerships with coworking spaces and small-business accountants. That gives contribution-margin payback of under 1.5\(\times\) per case at the \(\$29\) flat price and roughly 3\(\times\) per case once Phase 2 take-rate pricing is live. Per-case payback matters more than LTV here because small claims is episodic.

{\itshape Breakeven and margin risks.} Fixed monthly burn at launch is roughly \(\$8\)--\(\$12\)k (Cloud Run minimums, Vertex AI baseline, part-time-equivalent founder time), so breakeven sits near 400 cases/month at flat \(\$29\) and roughly 150 cases/month once Phase 2 is live---the take-rate transition is what makes the business compound rather than just survive. The margin killers are concentrated and known: huge evidence sets that explode token cost beyond the \(\$0.07\) baseline, refund and chargeback exposure if packet quality slips, Stripe fees of about 4\% on a \(\$29\) ticket, high-touch support when service-of-process bounces, and jurisdiction-expansion engineering when the wedge moves outside NYC.

{\bfseries Why these technical choices.} The system mirrors the real filing workflow. The Next.js wizard collects court-required facts; the Extractor turns messy evidence into structure; the DefendantResolver checks the NY Department of State registry; and the JurisdictionChecker uses deterministic tools plus a small Chroma RAG corpus to validate cap, statute of limitations, venue, and citations. Structured Pydantic outputs make the PDF renderer reliable. Gemini 2.5 Flash fits because it supports multimodal, long-context evidence review while keeping token cost far below the \(\$29\) price.

{\bfseries Risks and mitigations.}\\[-1pt]
\textbf{1. Bar and UPL risk:} ClaimReady is form preparation and educational guidance, not legal advice; every output carries a disclaimer, and complex cases are flagged for attorney review.\\[-1pt]
\textbf{2. Data dependency:} NY DOS is the source of truth for defendant lookup, mitigated by caching, manual verification fallback, and a per-state registry plug-in path.\\[-1pt]
\textbf{3. Model and quality risk:} Pydantic outputs, deterministic jurisdiction tools, RAG citations, and a packet-review checkpoint bound the LLM's surface area.\\[-1pt]
\textbf{4. Refund and chargeback risk:} the \(\$29\) ticket caps exposure, and failure modes are concrete enough to support a ``court-ready or your money back'' guarantee.

{\bfseries Go-to-market.} Three channels at launch with a partnership path layered on top.\\[-1pt]
\textbf{1.} SEO on the ``how do I sue an LLC in NY'' long-tail, with borough-specific landing pages and anonymized case-study posts.\\[-1pt]
\textbf{2.} Community presence in Reddit communities for freelancers and personal finance, freelancer Twitter, and the Freelancers Union mailing list.\\[-1pt]
\textbf{3.} Local partnerships with coworking spaces, small-business accountants, and freelance-friendly co-ops, where ClaimReady is offered at a member-benefit discount.\\[-1pt]
\textbf{4.} The Phase 3 path: B2B distribution inside freelancer-facing platforms where the partner pays acquisition and ClaimReady captures the engine economics---unlocked once consumer-channel CAC and packet-quality metrics are proven.

{\bfseries MVP and wedge.} The MVP is deliberately narrow: NYC small-claims, unpaid-invoice disputes, and NY-registered business defendants. NYC is the wedge because it has dense freelancer demand, a generous \(\$10{,}000\) cap, and public DOS entity data. Winning this workflow gives ClaimReady real users, packet data, defendant patterns, and a reusable filing engine. Expansion then follows two paths: \textbf{claim types} such as security deposits, consumer-services contracts, unreturned property, and vendor disputes; and \textbf{jurisdictions} such as CA, TX, IL, and FL, where the same architecture can swap in local caps, venue rules, forms, and entity registries.
\vspace{-1pt}

{\bfseries Close.} ClaimReady wins by being deliberately small at launch: not ``AI lawyer for everything,'' but a cheap, focused filing assistant for one painful, common, economically underserved workflow, designed from day one as the wedge for a broader civil-claims platform.

\vspace{1pt}
{\scriptsize
Sources: \href{https://www.upwork.com/press/releases/upwork-study-finds-64-million-americans-freelanced-in-2023-adding-1-27-trillion-to-u-s-economy}{Upwork Freelance Forward 2023};
\href{https://nwu.org/survey-shows-that-62-of-ny-freelancers-experience-non-payment-need-freelance-isnt-free-law-statewide/}{NY freelancer nonpayment survey};
\href{https://nycourts.gov/courts/nyc/smallclaims/forms/instructionsforfiling.pdf}{NYC Small Claims filing instructions};
\href{https://www.clio.com/resources/legal-trends/compare-lawyer-rates/ny/}{Clio NY lawyer rates};
\href{https://cloud.google.com/vertex-ai/generative-ai/pricing}{Google Vertex AI pricing}.}

\end{document}
